\section{Compact quotients beyond Sol}\label{sec:catalog}
\subsection{Spherical space forms}
A round unit sphere has spin levels \(\pm(k+3/2)\), \(k=0,1,\ldots\). Let \(V_j=\mathrm{Sym}^j(\mathbb C^2)\). For a lifted finite group \(\widetilde G\subset SU(2)_L\times SU(2)_R\), restrict \(V_k\boxtimes V_{k+1}\) and \(V_{k+1}\boxtimes V_k\) to \(\widetilde G\) and take invariants. Their dimensions \(m_+(k),m_-(k)\) are the two multiplicities in the convention of Appendix~\ref{app:spheres}. This procedure commutes with the scalar shift. It gives, on the round maximal-expansion surface of the minimal model, kernel dimensions \((2,0)\) for the left order-three quotient, \((0,2)\) for its mirror, \((0,0)\) for \(L(5,2)\) and \((2,2)\) for the cover, in positive/negative charge order. The magnitudes are fixed by \(q^2_{\max}=|Q_\Sigma|/(6V_0)\) with \(q\) the round radius and \(V_0\) unit-radius volume, hence \(|qQ_\Sigma/(4\vol)|=3/2\) there. The same lowest levels can thus give a one-sided contact, no contact or both-sided contact.

Projecting the spinor representation also gives the full eta invariant. This distinction matters for \(L(7,2)\): both first shells can be missing although \(\eta\ne0\). The 15 lift entries in Appendix~\ref{app:spheres} keep one orientation throughout. The same three-dimensional table applies wherever this round slice occurs, independent of eventual four-dimensional boundary conditions.

\subsection{Nil: chirality survives a collision of thresholds}
On a Heisenberg quotient let the central generator be \(\gamma\) and the horizontal commutator \([x,y]=\gamma^r\), \(r\ge1\). An automorphism with horizontal integral matrix \(P\) and central sign \(\epsilon\) satisfies \(\gamma^{r\det P}=\gamma^{r\epsilon}\), hence \(\det P=\epsilon\). The total degree is \(+1\). Every quotient in this family is therefore chiral. Spin characters \((-1)^{2\delta_i}\) satisfy \(2r\delta_3=0\bmod2\); this leaves four structures for odd \(r\), eight for even \(r\).

For \(g=a^2(dx^2+dy^2)+c^2(dz-x\,dy)^2\), set \(a=1,c=t\) and measure the nonzero-centre sector. With \(k_*=r\) for central-periodic spin and \(k_*=r/2\) otherwise, put \(\nu_*=k_*/(2\pi)\) and \(R_*=\sqrt{(\nu_*/t)^2+2\nu_*}\). On the tested leaves \(t=1/2,1,2\), \(R_*>t/4\), and comparison with the lowest unpaired level gives
\begin{equation}
 \frac{q_+}{6V_0}=\frac{2t}{3}(R_*-t/4),\qquad
 \frac{q_-}{6V_0}=\frac{2t}{3}(\nu_*/t+t/4).
 \label{eq:nilmin}
\end{equation}
At \(t=1\) the two magnitudes differ for all 24 spin structures with \(1\le r\le4\). For the basic lattice the numbers are \(0.224138873\) and \(0.272769962\). Setting \(t=2\sqrt{\nu_*}\) makes the two equal and keeps the stated ground branches lowest. Chirality has remained fixed. Adding the centre-zero torus sector can introduce a nearer pair; e.g.\ at \(r=3,4\), \(t=2\) and trivial spin it gives equal full magnitudes \(q_\pm/(6V_0)=2/3\). Appendix~\ref{app:nil} derives the three spectral towers and the eta correction.

\subsection{Flat space forms}
The six orientable compact flat three-manifold types \(G_1,\dots,G_6\) have respectively \(8,8,2,4,2,4\) spin structures. There are eight with nonzero eta, two each on the half-, third-, quarter- and sixth-turn spaces, which we label \(G_2\) through \(G_5\). The half-turn manifold admits orientation reversal. Its screw generator \(S\) is sent under central inversion to \(Z^{-1}S\) with \(Z\) the axial translation. The lift of \(Z\) can be negative, so a lift with a spectrally asymmetric axial tower is carried to the other lift. Conversely on the chiral quarter-turn manifold a horizontally twisted lift removes the axial tower; the remaining modes occur in signed pairs.

For horizontal-periodic spin on a screw manifold of order \(m\), the asymmetric tower has eigenvalues \(2\pi(mj+c_*)/L\), each twice, with \(0<c_*<m\). Thus \(\eta=2(1-2c_*/m)\). The nearest nonaxial orbit supplies both signs equally at \(2\pi\sqrt{\rho^2+z_0^2/L^2}\), where \(\rho\) is the shortest nonzero horizontal covector length and \(z_0\) the least axial frequency in translation units. This level can precede both axial levels. On \(G_5\), negative screw lift, equality persists through \(L=3/\sqrt2\), even though \(\eta=1/3\) throughout. Lengthening the axis past this value splits the pair without changing eta. Appendix~\ref{app:flat} gives the spin and normalization dictionary.

\subsection{Product and Seifert geometries}
In \(S^2\times S^1\) the angular modes pair at each circle momentum. The quotient \(\mathbb{RP}^3\#\mathbb{RP}^3\) combines the antipodal map of \(S^2\) with two circle reflections. A periodic spinor retains a zero-momentum tower that alternates sign with angular momentum. This gives ordinary eta \(+1/2,-1/2\) for the two periodic lifts, zero for the two antiperiodic lifts. An orientation reversal exchanges the periodic lifts. Appendix~\ref{app:product} gives the spectrum and its continuation.

As a comparison of conventions, the Brieskorn links \(M(2,3,5)\), \(M(2,3,6)\), \(M(2,3,7)\) realize three different geometries. Under the metrics and orientation specified in Appendix~\ref{app:product}, the ordinary spin eta values at unit aspect are respectively \(1079/720\), \(43/24\) or \(-5/24\) according to spin, and \(-1853/1008\). These are controls on conventions and metric dependence. A general hyperbolic-base Seifert quotient requires additional base-spectrum data to locate its first threshold.
